% Lección: el teorema de Bolzano (castellano).
%
% Teorema, demostración, corolario y ejemplo. La demostración va en
% `notesonly`: en las diapositivas se enuncia y en los apuntes se demuestra.

\begin{frame}
\didactatitle{El teorema de Bolzano}

\begin{theorem}[Bolzano]
Si $f$ es continua en $[a, b]$ y $f(a)$ y $f(b)$ tienen signos opuestos,
existe $c \in (a, b)$ tal que $f(c) = 0$.
\end{theorem}

% \bymedium{corto}{largo}: la primera versión en las diapositivas y la
% segunda en los apuntes. Casi siempre es mejor que esconder la frase.
\bymedium{%
  Una función continua que empieza por debajo del eje y acaba por encima
  tiene que cruzarlo.
}{%
  Una función continua que empieza por debajo del eje y acaba por encima
  tiene que cruzarlo. Lo que el teorema añade a la intuición es que eso no
  depende del dibujo, sino de que la recta real no tiene huecos: el axioma
  del supremo.
}

\begin{notesonly}
\begin{proof}
Supongamos $f(a) < 0 < f(b)$; el otro caso se reduce a éste cambiando $f$ por
$-f$. El conjunto $A = \{x \in [a, b] : f(x) < 0\}$ contiene a $a$ y está
acotado por $b$, así que tiene supremo $c \in [a, b]$. Veamos que $f(c) = 0$.

Si fuera $f(c) \neq 0$, tomando $\varepsilon = |f(c)|$ en la definición de
continuidad habría un $\delta > 0$ tal que $f(x)$ tiene el mismo signo que
$f(c)$ para todo $x \in [a, b]$ con $|x - c| < \delta$. Si $f(c) < 0$,
entonces $c < b$ y habría puntos de $A$ mayores que $c$. Si $f(c) > 0$,
entonces $c > a$ y ningún punto de $(c - \delta, c]$ estaría en $A$, así que
$c - \delta$ sería una cota superior de $A$ menor que $c$. Las dos cosas
contradicen que $c$ sea el supremo.

Por tanto $f(c) = 0$, y $c \neq a, b$ porque ni $f(a)$ ni $f(b)$ se anulan.
\end{proof}
\end{notesonly}

\begin{teaching}[Para discutir]
¿Qué falla si $f$ no es continua? (El salto $H(x) - \tfrac12$ en $[-1, 1]$.)
¿Y si solo
existieran los racionales? ($x^2 - 2$ en $[0, 2]$ cambia de signo y no tiene
ninguna raíz racional.)
\end{teaching}
\end{frame}

\begin{frame}
\didactatitle{Valores intermedios}

\begin{corollary}[Valores intermedios]
Si $f$ es continua en $[a, b]$, toma todos los valores comprendidos entre
$f(a)$ y $f(b)$.
\end{corollary}

\onlynotes{%
  Si $k$ está estrictamente entre $f(a)$ y $f(b)$, basta aplicar el teorema de
  Bolzano a $g(x) = f(x) - k$, que es continua y cambia de signo en $[a, b]$.
}

\dpause

\begin{example}
La ecuación $x^3 + x - 1 = 0$ tiene una solución en $(0, 1)$: el polinomio es
continuo, vale $-1$ en $0$ y $1$ en $1$.
\onlynotes{%
  Repitiendo el argumento en la mitad del intervalo donde cambia el signo, y
  luego en la mitad de esa mitad, se obtiene el \keyterm{método de
  bisección}: cada paso divide por dos la longitud del intervalo en que está
  la raíz.
}
\end{example}
\end{frame}
